% Fuente: Generalización del lagrangiano. Mismos cuatro momentos.
\begin{frame}[label=nav-frame-073]{Diccionario de los momentos generalizados}
\TablaSetup
\begin{tabular}{@{}p{.20\textwidth}p{.32\textwidth}p{.41\textwidth}@{}}
\TablaCabecera{Argumento} & \TablaCabecera{Momento} & \TablaCabecera{Qué variación acompaña} \\
\TablaLinea
$R_{abcd}$ & $P^{abcd}=\dfrac{\partial L}{\partial R_{abcd}}$ & $P^{abcd}\,\delta R_{abcd}$\newline Variación de la curvatura. \\[3mm]
$g^{ab}$ & $P_{ab}=\dfrac{\partial L}{\partial g^{ab}}$ & $P_{ab}\,\delta g^{ab}$\newline Variación de la métrica. \\[3mm]
$\nabla_a\phi$ & $\Pi^a=\dfrac{\partial L}{\partial(\nabla_a\phi)}$ & $\Pi^a\,\delta(\nabla_a\phi)$\newline Variación del gradiente. \\[3mm]
$\phi$ & $F_\phi=\dfrac{\partial L}{\partial\phi}$ & $F_\phi\,\delta\phi$\newline Dependencia explícita en $\phi$. \\
\TablaLinea
\end{tabular}\par
\vspace{3mm}
\footnotesize
Cada derivada mantiene fijos los demás argumentos de $L$.
Aquí $P_{ab}$ deriva respecto de $g^{ab}$, a diferencia del
$P^{ab}$ inicial, definido respecto de $g_{ab}$.
\end{frame}
